% Generated draft: requires mathematical and source review.
\documentclass[11pt]{article}
\usepackage[margin=25mm]{geometry}
\usepackage{fontspec}
\setmainfont{Latin Modern Roman}
\usepackage{amsmath,amssymb,mathtools,unicode-math}
\setmathfont{Latin Modern Math}
\usepackage{enumitem,needspace,xurl,hyperref}
\setlength{\parindent}{0pt}
\setlength{\parskip}{5pt}
\setlength{\emergencystretch}{3em}
\newcommand{\N}{\mathbb{N}}
\newcommand{\Z}{\mathbb{Z}}
\newcommand{\Q}{\mathbb{Q}}
\newcommand{\R}{\mathbb{R}}
\newcommand{\C}{\mathbb{C}}
\newcommand{\F}{\mathbb{F}}
\newcommand{\A}{\mathbb{A}}
\newcommand{\PP}{\mathbb{P}}
\newcommand{\E}{\mathbb{E}}
\newcommand{\eps}{\varepsilon}
\newcommand{\dd}{\,\mathrm{d}}
\begin{document}
\section*{UnsolvedMath ID: 11000040}
\textbf{Source problem code:} catalogue entry
\subsection{Definitions and mathematical statement}
Let $\mathcal T_g$ be the Teichm\"uller space of marked closed genus-$g$ Riemann surfaces, where $g\ge2$, and let $\operatorname{Mod}_g$ act by changing markings. Let $h$ be a $\operatorname{Mod}_g$-invariant Riemannian metric, or a Finsler metric with the source's unspecified weak regularity conditions. Finite covolume means that the moduli orbifold $\operatorname{Mod}_g\backslash\mathcal T_g$ has finite $h$-volume, with a suitable volume convention needed in the Finsler case. Write $\Gamma_g$ for the image of the mapping class action and $\operatorname{Isom}(\mathcal T_g,h)$ for the full isometry group with its usual compact-open topology.

\textbf{Conjecture 4.1 (inhomogeneity of all metrics).} For every such $h$, the isometry group is discrete and the mapping class action has uniformly bounded finite index in it:
\[
 [\operatorname{Isom}(\mathcal T_g,h):\Gamma_g]\le C(g).
\]
The source phrases the index as $C=C(g)$; the display uses the natural uniform-bound interpretation. If literal equality independent of $h$ is intended, that is a distinct stronger statement requiring clarification. The claim concerns all metrics under the given invariance and covolume hypotheses, not only familiar natural metrics. Completeness is not an explicit hypothesis in the extracted statement and has not been inserted.

The source writes $\operatorname{Mod}_g$ itself as an isometry subgroup. The image notation matters in genus two, where the closed-surface mapping class action has a central kernel; the literal abstract subgroup wording needs an effectiveness qualification.

The source proposes the following extensions and refinements:
\begin{enumerate}
\item Replace $\operatorname{Mod}_g$ by any finite-index subgroup $H$, using an $H$-invariant finite-covolume metric, and seek discreteness and a bound depending only on $g$ and $[\operatorname{Mod}_g:H]$ for the index of the image of $H$ in the full isometry group.
\item After this, consider substantially smaller groups such as the Torelli group $\mathcal I_g$, at least seeking discreteness when the group is ``sufficiently large.'' No largeness condition or complete metric/covolume hypotheses for this exploratory variant are supplied.
\item Seek the sharper possibility that the index constant can be taken to be $1$ or $2$. The source suggests studying $\mathcal T_g/\Lambda$ for discrete isometry groups $\Lambda$ properly containing the mapping class image, as orbifolds finitely covered by moduli space $\mathcal M_g$.
\end{enumerate}
It cites Royden's theorem for the Teichm\"uller-metric case, reporting index $2$, and suggests extending [FW1] from closed to finite-volume settings as evidence. The discussion attributes inhomogeneity to global moduli-space topology rather than solely to the fine unit-ball geometry in the spaces $Q(X)$ of holomorphic quadratic differentials.
\subsection{Short English statement}
Every mapping-class-invariant finite-covolume metric on Teichm\"uller space is conjectured to have a discrete isometry group containing the mapping class image with uniformly bounded finite index. The source also asks for finite-index-subgroup versions and a possible index bound of two.
\subsection{Sources}
\begin{itemize}
\item Local UnsolvedMath registry, permanent ID 11000040, source code catalogue entry; no individual proposer or proposal year supplied.
\item Farb, edited volume \emph{Problems on Mapping Class Groups and Related Topics}, Conjecture 4.1, PDF page 37; \url{https://www.math.uchicago.edu/~farb/papers/mcgbook.pdf}. The registry records PDF-font-structure extraction and access on 2026-07-29; its inventory has 329 labeled headings distinguished by source page.
\item The excerpt cites Royden's theorem and source label [FW1], without full bibliographic entries. It reports the Royden index-$2$ case and a proposed closed-to-finite-volume extension of [FW1], not an established extension.
\end{itemize}
\end{document}
